\SourceLeaf{374}%
von \(AO\) ein Stück \(OM=NL\) abschneidet und \(ML\) zieht, wenn man ferner in \(M\) an \(MA\) wieder einen Winkel von der Größe des Winkels \(A\) anträgt, dessen anderer Schenkel \(MC\) die Gerade \(AN\) im Punkt \(C\) zwischen \(A\) und \(L\) trifft, darauf von \(AM\) ein Stück \(MB=LC\) abschneidet und \(BC\) zieht und diese Konstruktion dann in ähnlicher Weise fortsetzt, sodass auf der Geraden \(OA\) die Punkte \(O,M,B,E,G,K\) usw. und auf der Geraden \(NA\) die Punkte \(N,L,C,D,F,H\) usw. liegen, so wird behauptet, dass die Stücke \(OM,MB,BE,EG,GK\) usw. oder die ihnen jeweils gleichen Stücke \(NL,LC,CD,DF,FH\) usw. eine abnehmende Progression bilden.

Der Verfasser versucht, diesen Lehrsatz durch Widerspruch zu beweisen. Er zählt die anderen möglichen Fälle auf, die sich bei einem falschen Lehrsatz ergeben könnten, und versucht zu zeigen, dass jeder von ihnen ausgeschlossen ist. Er behauptet nämlich, dass unter dieser Voraussetzung einer der folgenden fünf Fälle eintreten müsste. Die aufeinanderfolgenden Stücke, von \(OM\) an gerechnet, wären
\par\noindent\hspace*{\parindent}1) alle einander gleich, oder
\par\noindent\hspace*{\parindent}2) jedes folgende größer als das vorhergehende, oder
\par\noindent\hspace*{\parindent}3) einige gleich groß und das darauf folgende größer oder kleiner, oder
\par\noindent\hspace*{\parindent}4) einige aufeinanderfolgende Stücke nähmen fortschreitend ab und die darauf folgenden fortschreitend zu oder
\par\noindent\hspace*{\parindent}5) sie würden abwechselnd größer und kleiner.
\par\noindent In dieser Aufzählung fehlt der mögliche Fall, dass die Stücke zunächst fortschreitend zu- und dann fortschreitend abnehmen. Nach der eigenen Überzeugung des Rezensenten (deren tiefer liegende Gründe hier aber nicht angeführt werden können) wäre es gerade die Hauptsache und die eigentliche Lösung des Gordischen Knotens, diesen Fall zu erledigen. Allerdings kann man zugeben, dass diese Auslassung hier insofern wenig ausmacht, als die Beweisweise des Verfassers, mit der er den dritten Fall ausschließen will, sich von selbst auch auf diesen Fall erstrecken ließe, falls sie zulässig wäre. Doch gerade diesem angeblichen Beweis dafür, dass der dritte Fall ausgeschlossen ist, können wir keine Gültigkeit zugestehen. Der Verfasser stellt die Sache folgendermaßen dar. Nimmt man etwa im dritten Fall an, dass die beiden ersten Stücke gleich sind und das dritte größer ist, so wäre \(DC\) also größer als \(CL\). Nun ist aber \(AML\) ebenfalls ein gleichschenkliges Dreieck, das dieselbe Grundbedingung erfüllt wie das ursprüngliche Dreieck \(AON\). Wenn also jener dritte Fall mit seinem angenommenen Unterfall gelten würde,
